\documentclass[border=12pt]{standalone}
\usepackage{ctex}
\usepackage{amsmath,amssymb}
\usepackage{xcolor}
\setlength{\parindent}{0pt}
\setlength{\parskip}{6pt}

\definecolor{titleblue}{RGB}{24,62,143}
\definecolor{lineblue}{RGB}{74,118,196}
\definecolor{bgblue}{RGB}{246,249,255}

\newcommand{\stage}[2]{%
  \vspace{2mm}
  \noindent\fcolorbox{lineblue}{bgblue}{%
    \begin{minipage}{0.95\linewidth}
      \textbf{#1}\\[2pt]
      #2
    \end{minipage}%
  }%
}

\begin{document}
\begin{minipage}{22.5cm}
{\LARGE\bfseries\color{titleblue} Actor-Critic（Advantage）算法流程}

\textbf{输入：}策略网络 $\pi_\theta$，价值网络 $V_\phi$，学习率 $\alpha_\theta,\alpha_\phi$，折扣因子 $\gamma$。\\
\textbf{初始化：}初始化参数 $\theta,\phi$（LLM 场景通常从 SFT 参数继续训练）。

\textbf{每个训练回合（episode）包含以下 5 个阶段：}

\stage{阶段 A：采样轨迹}{
使用当前策略与环境交互，得到轨迹
\[
\tau=(s_0,a_0,r_0,s_1,a_1,r_1,\ldots,s_T)
\]
}

\stage{阶段 B：逐时刻估计优势（TD 误差）}{
对每个时间步 $t$ 计算
\[
A_t=r_t+\gamma V_\phi(s_{t+1})-V_\phi(s_t)
\]
}

\stage{阶段 C：更新 Critic（价值网络）}{
令 TD 目标为
\[
y_t=r_t+\gamma V_\phi(s_{t+1})
\]
最小化损失
\[
L_{\text{critic}}=\bigl(V_\phi(s_t)-y_t\bigr)^2
\]
}

\stage{阶段 D：更新 Actor（策略网络）}{
沿策略梯度方向更新
\[
\theta\leftarrow\theta+\alpha_\theta\,\nabla_\theta\log\pi_\theta(a_t\mid s_t)\cdot A_t
\]
}

\stage{阶段 E：进入下一轮}{
重复阶段 A--D，直到策略收敛或达到最大迭代次数。
}
\end{minipage}
\end{document}
